\section{Encapsulation}
The module is intended for operation in harsh environments where exposure to water, oil, and other contaminants may occur. To ensure reliable operation and long-term durability, dedicated protective enclosures and encapsulation solutions have been designed for both the electronics and sensor assemblies. These protective measures shield the components from environmental exposure while maintaining the functionality and performance of the system.

For the PCB, a silicone potting compound from \textit{TermoPasty} was selected as the encapsulation material \cite{termopasty_silicone_compound_011}. After curing, the compound forms a flexible rubber-like material, in contrast to epoxy-based potting compounds, which cure into a rigid solid. The flexible nature of the silicone encapsulation is advantageous for this application, as it allows the assembly to accommodate thermal expansion and contraction over the expected operating temperature range. This reduces mechanical stress on the PCB and components, thereby improving long-term reliability.

The compound has a mixing ratio of 100:8 of the silicone potting compound to the catalyzer. The two compound were measured in two different plastic cups before mixing. After a thorough mix the solution was poured over the PCB. The PCB was placed in a 3D printed box, working as a mold for the potting. 
A picture of the mixing process is shown in \cref{fig:silicone} and the cured solution is shown on \cref{fig:pcb_silicone}. 

\begin{figure}[H]
    \centering

    \begin{subfigure}{0.48\linewidth}
        \centering
        \includegraphics[width=\linewidth]{Pictures/Silicone_potting.jpeg}
        \caption{Preparation of potting silicone.}
        \label{fig:silicone}
    \end{subfigure}
    \hfill
    \begin{subfigure}{0.48\linewidth}
        \centering
        \includegraphics[width=1.22\linewidth]{Pictures/silicon_sideview_nobg.png}
        \caption{PCB after potting.}
        \label{fig:pcb_silicone}
    \end{subfigure}

    \caption{Silicone potting process and result.}
    \label{fig:potting_comparison}
\end{figure}

